% ============================================================================
% Shared preamble for the Plexus CRDT preprint stack
% Adapted from WebstormProjects/plato main.tex, with arxiv-style 2-column layout:
% title + abstract share the first page with content flowing into columns below.
%
% Typography-heavy packages (libertinus, microtype expansion, titlesec, listings)
% are loaded conditionally via \IfFileExists so the paper compiles on minimal
% TeX Live installs. The production build uses Docker (texlive:latest) which
% has all of these — see build.sh.
% ============================================================================

\documentclass[10pt,twocolumn]{article}

\usepackage[
    margin=0.75in,
    columnsep=0.3in,
    heightrounded
]{geometry}

\usepackage[square,sort&compress,super,comma,numbers]{natbib}
\usepackage{graphicx}
\usepackage{etoolbox}
\usepackage{booktabs}
\usepackage{tabularx}
\usepackage{amsmath}
\usepackage{amssymb}
\usepackage[utf8]{inputenc}
\usepackage{csquotes}
\usepackage{url}
\usepackage[T1]{fontenc}
\usepackage{textcomp}
\usepackage[normalem]{ulem}
\usepackage{setspace}
\usepackage{xcolor}

% Font: libertinus when available; otherwise default (Computer Modern).
% alphabeta (Greek support) only makes sense with libertinus — CM Greek
% encoding files are not always bundled in minimal TeX Live.
\IfFileExists{libertinus.sty}{%
    \usepackage{libertinus}%
    \IfFileExists{alphabeta.sty}{\usepackage{alphabeta}}{}%
}{}

% Microtype: expansion requires scalable fonts — only enable when libertinus present.
\IfFileExists{libertinus.sty}{%
    \usepackage[protrusion=true, expansion=true]{microtype}%
    % Protrusion tuning only meaningful when microtype is loaded.
    \SetProtrusion[ name = default ]{ encoding = {*} }{ "2014 = {0,0} }%
    \SetProtrusion{ encoding = {*}, family = {*} }{%
        `     = {400,0},%
        ''    = {0,400},%
        .     = {0,600},%
        ,     = {0,500},%
        :     = {0,400},%
        ;     = {0,400}%
    }%
    \SetTracking{encoding={*}, shape=sc}{40}%
}{}

% ============================================================================
% TYPOGRAPHY REFINEMENTS
% ============================================================================

\setstretch{1.1}

\widowpenalty=10000
\clubpenalty=10000
\raggedbottom

\hyphenpenalty=2000
\exhyphenpenalty=2000

\renewenvironment{quote}{%
    \list{}{%
        \leftmargin=1.5em%
        \rightmargin=1.5em%
        \parsep=0pt%
        \topsep=8pt plus 2pt minus 2pt%
    }%
    \item\relax\small%
}{%
    \endlist%
}

\IfFileExists{footmisc.sty}{\usepackage[bottom]{footmisc}}{}
\setlength{\footnotesep}{1.1\baselineskip}

\renewcommand{\arraystretch}{1.2}
\setlength{\tabcolsep}{6pt}

% ============================================================================
% SECTION STYLING (optional titlesec)
% ============================================================================

\IfFileExists{titlesec.sty}{%
    \usepackage{titlesec}%
    \titleformat{\section}%
        {\normalfont\large\bfseries}{\thesection.}{0.5em}{}%
    \titleformat{\subsection}%
        {\normalfont\normalsize\bfseries}{\thesubsection}{0.5em}{}%
    \titleformat{\subsubsection}%
        {\normalfont\small\bfseries\itshape}{}{0pt}{}%
    \titlespacing*{\section}{0pt}{12pt plus 2pt minus 2pt}{6pt}%
    \titlespacing*{\subsection}{0pt}{10pt plus 2pt minus 2pt}{4pt}%
    \titlespacing*{\subsubsection}{0pt}{8pt plus 1pt minus 1pt}{2pt}%
}{}

% ============================================================================
% CODE / PSEUDOCODE LISTINGS
% ============================================================================

\IfFileExists{listings.sty}{%
    \usepackage{listings}%
    \lstset{%
        basicstyle=\ttfamily\footnotesize,%
        breaklines=true,%
        breakatwhitespace=false,%
        columns=fullflexible,%
        keepspaces=true,%
        frame=single,%
        framesep=3pt,%
        framerule=0.4pt,%
        xleftmargin=0pt,%
        xrightmargin=0pt,%
        aboveskip=6pt,%
        belowskip=6pt,%
        showstringspaces=false,%
        numbers=none%
    }%
}{%
    % Fallback: lstlisting is not available; define a minimal alias mapping
    % \begin{lstlisting}...\end{lstlisting} to a verbatim environment.
    \usepackage{verbatim}%
    \let\lstlisting\verbatim%
    \let\endlstlisting\endverbatim%
}

% ============================================================================
% ARXIV-STYLE TITLE: title + abstract + content share the first page
% ============================================================================

\makeatletter
% Redefine \and to produce a horizontal gap rather than a tabular \cr.
% This lets \author{A \and B} render inline in our custom \maketitle.
\renewcommand{\and}{\hskip 2em\relax}

\renewcommand{\maketitle}{%
    \twocolumn[%
        \begin{@twocolumnfalse}%
        \begin{center}%
        {\Large\bfseries \@title \par}%
        \vspace{0.8em}%
        {\normalsize \@author \par}%
        \vspace{0.5em}%
        {\small \@date \par}%
        \end{center}%
        \vspace{1em}%
        \begin{abstract}%
        \@abstract%
        \end{abstract}%
        \vspace{1em}%
        \end{@twocolumnfalse}%
    ]%
    % Output any \thanks footnotes accumulated during author parsing.
    % With \twocolumnfalse these land on page 1, which is the intended
    % arxiv-style placement.
    \@thanks\let\@thanks\relax%
}

\long\def\@abstract{}
\long\def\plexusabstract#1{\long\gdef\@abstract{#1}}
\makeatother

% ============================================================================
% Hyperref (load late)
% ============================================================================

\usepackage{hyperref}
\hypersetup{
    colorlinks=true,
    linkcolor=blue!60!black,
    citecolor=blue!60!black,
    urlcolor=blue!60!black,
    pdfborder={0 0 0}
}
